\section{Related work}
\label{sec:related}

Several tools use \egraphs for rule inference~\citep{sat19, ruler}.
\citet{sat19} use enumerative
  synthesis to infer axioms for the CVC4 theorem prover.
Ruler~\citep{ruler} outperformed \citet{sat19}
  in various domains.
  In this paper,
  we show that we can outperform Ruler in terms of both
  scalability and generality of domains (\autoref{sec:eval}).

Similarly, theory exploration is a well-studied topic
 focusing on eagerly synthesizing
 lemmas that may be useful for
 verification tasks.
A recent tool in this space is
  TheSy~\citep{cav21}, which performs 
  inductive theory exploration using
  \eqsat and symbolic values
  to efficiently filter candidate conjectures.
TheSy's key insight is to leverage congruence closure
 to implement an induction prover within the
 \eqsat framework.
Other, similar tools for theory exploration
 use random testing to find
 potential candidates~\citep{hipspec, quickspec}.
IsaCoSy~\citep{isacosy} synthesizes inductive
  theorems for the Isabelle theorem prover~\citep{isabelle}.
To keep the search space of terms tractable,
  IsaCoSy selectively enumerates only terms that
  are not reducible from existing rules.
A similar technique is used by \citet{seplogic}
  in lemma synthesis for proving
  entailments with separation logic.
We believe the abstractions provided by \enumo for
  guided search and ruleset manipulation
  can be used to scale lemma synthesis
  in these tools.
In future work, we would like to express the
inductive prover from  TheSy
 in \enumo.

Many custom tools have been proposed
  for synthesizing rewrite rules in
  specific domains.
\citet{queso} proposed a tool that
  synthesizes rewrite rules for
  quantum circuit optimization.
\citet{taso19} developed a tool synthesizing
  graph substitutions for deep neural networks.
RuleSy~\citep{rulesy} uses a combination of
synthesis and specification mining to find
proof rules representing the mined
specification.
\citet{wetune} presented a tool for
  discovering and verifying query rewrite rules.
In contrast,
  \enumo's DSL-based approach is
  not specialized to any particular domain;
  our evaluation in \autoref{sec:eval}
  shows that \enumo works across diverse domains.
Prior work has used machine learning to assist in
  rewrite rule inference~\citep{oopsla22-logic, swapper};
in particular, \citet{swapper} supports some
  forms of conditional rules.
More recently, large language models have shown promise
  for synthesis tasks~\citep{hysynth, asap, shraddha};
  in \autoref{sec:case-study-llm}, we show that \enumo's
  operators make it easy to incorporate LLMs into
  theory exploration, where they complement guided search.
This paper shows that using \enumo's novel term enumeration
  primitives, rule inference scales
  to support grammars that have conditional
  operators; however,
  full support for conditional rule inference
  is left for future work.

Finally, several tools have focused on automatically inferring
  peephole optimizations~\citep{bansal, alive, fraser, optgen}, and
  instruction selection~\citep{buchwald}.
Two major challenges with these optimizations
  are the presence of side conditions
  and their large grammars.
This paper shows that
  with \enumo's guided enumeration strategy,
  it is possible to find rewrite rules with 
  side conditions.
  We also show that it is possible to
  scale to
  large grammars, like that of Halide~\citep{halide}.
We will continue to explore
  better support for more general conditional
  rewrite rule inference,
  and we are excited to use \enumo
  to infer more optimizations for frameworks like LLVM.

The \enumo DSL is designed to facilitate
  efficient term enumeration given a grammar.
Effective enumeration has been explored in many other contexts like
  relational algebra, sorting algorithms, testing, and generating well-typed
  lambda terms~\citep{alice, icfp16, feat, flajolet, welltyped}.
In the most closely related work, \citet{feat}
  propose \textit{Feat}, a Haskell library for composing enumerations.
They use a lazy mechanism (\textit{functional enumeration})
  to scale enumeration and leverage memoization to efficiently index into a stream
  of enumerated terms.
In a previous prototype of \enumo, we explored a similar mechanism,
  but we found that in the context of rewrite rule inference,
  enumeration time is never the bottleneck---working
  with a large \egraph is.
Therefore, in our final prototype, we use a simpler method for materializing
  a workload into a concrete set of terms.
A similarity between the \enumo DSL and
  the Feat library is the idea of
  composable workloads.
\citet{feat} define a set of combinators that
  allow them to compose smaller
  enumerations effectively.
As \autoref{sec:slide} showed,
  \enumo workloads
  can be composed using
  a set of operators (\Plug, \Filter, \Union), some of which
  are similar to Feat's (e.g., \textit{union}ing two workloads or enumerations).
A key feature of \enumo is that it
  only evaluates a workload when
  converting it to an \egraph (as shown in \autoref{sec:slide})---this
  allows \enumo to leverage a unique set of operators like
  \C{plug} and \C{filter} to
  optimize a workload before it is evaluated
  and converted to an \egraph.

